%% kapa: exam=1 round=31 year=2020 period=2 form=A subject=accounting subject_ko=회계학 seq=7 number=47
%% source: Q-Net 감정평가사 기출문제 (cst00309) artlSeq=5207219 file=extracted/1008/5207219/A형/A형/2020년도 제31회 감정평가사 1차 시험_2교시_A형.pdf
%% conversion: kordoc 4.24.51 -> tools/gen1.py (auto); review=proofread-sonnet; flags=none
\begin{question}{47}
(주)감평은 20x1년 초 환경설비(취득원가 ￦5,000,000, 내용연수 5년, 잔존가치 ￦0, 정액법 상각)를 취득하였다. 동 환경설비는 관계법령에 의하여 내용연수가 종료되면 원상 복구해야 하며, 이러한 복구의무는 충당부채의 인식요건을 충족한다. (주)감평은 취득시점에 내용연수 종료 후 복구원가로 지출될 금액을 ￦200,000으로 추정하였으며, 현재가치계산에 사용될 적절한 할인율은 연 10\%로 내용연수 종료시점까지 변동이 없을 것으로 예상하였다. 하지만 (주)감평은 20x2년 초 환경설비의 내용연수 종료 후 복구원가로 지출될 금액이 ￦200,000에서 ￦300,000으로 증가할 것으로 예상하였으며, 현재가치계산에 사용될 할인율도 연 10\%에서 연 12\%로 수정하였다. (주)감평이 환경설비와 관련된 비용을 자본화하지 않는다고 할 때, 동 환경설비와 관련하여 20x2년도 포괄손익계산서에 인식할 비용은? (단, (주)감평은 모든 유형자산에 대하여 원가모형을 적용하고 있으며, 계산금액은 소수점 첫째자리에서 반올림하고, 단수차이로 인한 오차가 있으면 가장 근사치를 선택한다.)
\par{\small\centering\begin{tabular}{|c|c|c|}\hline
기간 & \makecell{단일금액 ￦1의 현재가치\\(할인율=10\%)} & \makecell{단일금액 ￦1의 현재가치\\(할인율=12\%)} \\ \hline
4 & 0.6830 & 0.6355 \\ \hline
5 & 0.6209 & 0.5674 \\ \hline
\end{tabular}\par}
\begin{choices}
  \item ￦1,024,837
  \item ￦1,037,254
  \item ￦1,038,350
  \item ￦1,047,716
  \item ￦1,061,227
\end{choices}
\end{question}
\answer{5}
